% !TEX root = exercices_s4.tex
% ============================================================================
% HEC Montréal MBA -- ECON50803: Macroeconomic Environment
% Shared Preamble for Exercise Handouts
% ============================================================================

% --- Document class ---------------------------------------------------------
\documentclass[11pt, letterpaper]{article}
\usepackage[margin=2.5cm]{geometry}

% --- Fonts ------------------------------------------------------------------
\usepackage[sfdefault,light]{FiraSans}
\usepackage{FiraMono}
\usepackage[T1]{fontenc}

% --- Colors (same as slides) ------------------------------------------------
\usepackage{xcolor}
\definecolor{HECnavy}{HTML}{002855}
\definecolor{HECgreen}{HTML}{26d07c}
\definecolor{HECcoral}{HTML}{ff585d}
\definecolor{HECtext}{HTML}{2D3748}
\definecolor{HECgray}{HTML}{94A3B8}
\definecolor{HEClightgray}{HTML}{F1F5F9}
\definecolor{HEClightblue}{HTML}{E8EDF3}

% --- Packages ---------------------------------------------------------------
\usepackage{amsmath, graphicx, tikz, booktabs, tabularx, hyperref}
\usepackage[inline]{enumitem}
\usepackage{etoolbox}
\usepackage{comment}
\usepackage{needspace}
\usepackage{tcolorbox}
\tcbuselibrary{skins, breakable}
\usepackage{titlesec}
\usepackage{fancyhdr}
\usepackage{lastpage}

% --- Paths ------------------------------------------------------------------
\graphicspath{{../Slides/Figures/}}

% --- Hyperlinks -------------------------------------------------------------
\hypersetup{
   colorlinks,
   linkcolor=HECnavy,
   urlcolor=HECnavy,
   citecolor=HECtext
}

% --- Answer toggle ----------------------------------------------------------
\newbool{answers}
\ifdefined\showanswers
   \booltrue{answers}
\else
   \boolfalse{answers}
\fi

% Answer environment: shown as green box when answers=true, excluded otherwise
\ifbool{answers}{%
   \newtcolorbox{reponse}[1][Réponse]{%
      colback=HECgreen!5,
      colframe=HECgreen,
      boxrule=0pt,
      leftrule=3pt,
      arc=4pt,
      title={\textcolor{HECgreen}{\bfseries #1}},
      coltitle=HECgreen,
      fonttitle=\bfseries,
      attach title to upper={\\[4pt]},
      left=8pt, right=8pt, top=4pt, bottom=4pt,
      before={\par\nopagebreak\medskip\nopagebreak}
   }%
}{%
   \excludecomment{reponse}%
}

% --- Custom box environments (adapted from slides) --------------------------

% Key Insight (green accent)
\newtcolorbox{keyinsight}[1][Point clé]{%
   colback=HECgreen!5,
   colframe=HECgreen,
   boxrule=0pt,
   leftrule=3pt,
   arc=4pt,
   title={\textcolor{HECgreen}{\bfseries #1}},
   coltitle=HECgreen,
   fonttitle=\bfseries,
   attach title to upper={\\[4pt]},
   left=8pt, right=8pt, top=4pt, bottom=4pt,
   breakable
}

% Business Implication (navy accent)
\newtcolorbox{bizimplication}[1][Implication d'affaires]{%
   colback=HECnavy!5,
   colframe=HECnavy,
   boxrule=0pt,
   leftrule=3pt,
   arc=4pt,
   title={\textcolor{HECnavy}{\bfseries #1}},
   coltitle=HECnavy,
   fonttitle=\bfseries,
   attach title to upper={\\[4pt]},
   left=8pt, right=8pt, top=4pt, bottom=4pt,
   breakable
}

% Warning / Important (coral accent)
\newtcolorbox{warning}[1][Important]{%
   colback=HECcoral!5,
   colframe=HECcoral,
   boxrule=0pt,
   leftrule=3pt,
   arc=4pt,
   title={\textcolor{HECcoral}{\bfseries #1}},
   coltitle=HECcoral,
   fonttitle=\bfseries,
   attach title to upper={\\[4pt]},
   left=8pt, right=8pt, top=4pt, bottom=4pt,
   breakable
}

% Definition (gray background)
\newtcolorbox{defbox}[1][Définition]{%
   colback=HEClightgray,
   colframe=HEClightgray,
   boxrule=0pt,
   leftrule=3pt,
   arc=4pt,
   left=8pt, right=8pt, top=4pt, bottom=4pt,
   title={\textcolor{HECnavy}{\bfseries #1}},
   coltitle=HECnavy,
   fonttitle=\bfseries,
   attach title to upper={\\[4pt]},
   breakable
}

% Exercise (coral accent)
\newtcolorbox{exercisebox}[1][Exercice]{%
   colback=HECcoral!5,
   colframe=HECcoral,
   boxrule=0pt,
   leftrule=3pt,
   arc=4pt,
   title={\textcolor{HECcoral}{\bfseries #1}},
   coltitle=HECcoral,
   fonttitle=\bfseries,
   attach title to upper={\\[4pt]},
   left=8pt, right=8pt, top=4pt, bottom=4pt,
   breakable
}

% --- Section formatting (titlesec) ------------------------------------------
\titleformat{\section}
   {\clearpage\Large\bfseries\color{HECnavy}}
   {\thesection.}{0.5em}{}
   [\vspace{2pt}{\color{HECgreen}\rule{2.5cm}{1.5pt}}]
\titlespacing*{\section}{0pt}{1.5em}{0.8em}

\titleformat{\subsection}
   {\large\bfseries\color{HECnavy}}
   {\thesubsection.}{0.5em}{}
\titlespacing*{\subsection}{0pt}{1em}{0.5em}

% --- Header/Footer (fancyhdr) ----------------------------------------------
\pagestyle{fancy}
\fancyhf{}
\renewcommand{\headrulewidth}{0pt}
\renewcommand{\footrulewidth}{0.4pt}
\renewcommand{\footrule}{\hfill{\color{HECgray}\rule{0.9\textwidth}{0.4pt}}\hfill\vspace{4pt}}
\fancyfoot[L]{\small\textcolor{HECgray}{ECON 50803\enspace\textbar\enspace Environnement macroéconomique}}
\fancyfoot[R]{\small\textcolor{HECgray}{\thepage\enspace/\enspace\pageref{LastPage}}}

% --- Enumitem defaults ------------------------------------------------------
\setlist[itemize]{leftmargin=1.5em, itemsep=3pt}
\setlist[enumerate]{leftmargin=1.5em, itemsep=3pt}

% --- MC question helpers ----------------------------------------------------
\newcounter{qcmcounter}
\newcommand{\qcm}[1]{%
   \stepcounter{qcmcounter}%
   \par\vspace{1.2em}%
   \needspace{12\baselineskip}%
   \noindent\textbf{\textcolor{HECnavy}{\theqcmcounter.}}\enspace #1%
   \vspace{0.3em}%
   \nopagebreak[4]%
}

\newcommand{\choix}[1]{%
   \nopagebreak[4]%
   \begin{enumerate}[label=\textcolor{HECgray}{(\alph*)}, leftmargin=2em, itemsep=1pt, topsep=2pt]
      #1
   \end{enumerate}
}

% --- Short answer counter ---------------------------------------------------
\newcounter{shortcounter}
\newcommand{\shortq}[1]{%
   \stepcounter{shortcounter}%
   \par\vspace{1.5em}%
   \needspace{4\baselineskip}%
   \noindent\textbf{\textcolor{HECnavy}{Question \theshortcounter.}}\enspace #1%
   \vspace{0.3em}%
}

% --- VFI (Vrai/Faux/Incertain) helper --------------------------------------
\newcommand{\vfi}[1]{%
   \stepcounter{shortcounter}%
   \par\vspace{1.5em}%
   \needspace{10\baselineskip}%
   \noindent\textbf{\textcolor{HECnavy}{Question \theshortcounter.}}\enspace\textit{Vrai, faux ou incertain?} #1%
   \vspace{0.3em}%
}

% --- Title page command -----------------------------------------------------
\newcommand{\exercisetitlepage}[2]{%
   % #1 = session number, #2 = session title
   \thispagestyle{empty}
   \begin{tikzpicture}[remember picture, overlay]
      % Navy sidebar
      \fill[HECnavy] (current page.north west) rectangle
         ([xshift=0.8cm]current page.south west);
      % Green accent line
      \fill[HECgreen] ([xshift=0.8cm]current page.north west) rectangle
         ([xshift=0.85cm]current page.south west);
   \end{tikzpicture}
   \vspace*{2cm}
   \hspace{0.8cm}%
   \begin{minipage}{0.85\textwidth}
      {\fontsize{28}{34}\selectfont\bfseries\textcolor{HECnavy}{Exercices --- Séance #1}}
      \\[14pt]
      {\Large\textcolor{HECnavy}{#2}}
      \\[14pt]
      {\color{HECgreen}\rule{3cm}{2pt}}
      \\[14pt]
      {\large\textcolor{HECgray}{ECON 50803 --- Environnement macroéconomique}}
      \\[8pt]
      {\large\textcolor{HECgray}{Jean-Félix Brouillette}}
      \\[20pt]
      \ifbool{answers}{%
         {\Large\bfseries\textcolor{HECgreen}{AVEC SOLUTIONS}}%
      }{%
         {\Large\bfseries\textcolor{HECgray}{SANS SOLUTIONS}}%
      }
   \end{minipage}
   \vfill
   \hspace{0.8cm}%
   {\small\textcolor{HECgray}{HEC Montr\'eal\enspace\textbar\enspace MBA}}
   \vspace{1cm}
   \newpage
}

% --- Mathematical commands (from slides) ------------------------------------
\newcommand{\pctchg}[1]{\%\Delta #1}


\begin{document}

\exercisetitlepage{4}{Cycles économiques et modèle AD-AS}

% ============================================================================
\section{Questions à choix multiples}
% ============================================================================

\setcounter{qcmcounter}{0}

\qcm{On observe simultanément une baisse du PIB réel, une hausse du chômage, une baisse de l'inflation et des prix du pétrole globalement stables. Quel diagnostic est le plus plausible dans le modèle AD-AS\,?}
\choix{
   \item Un choc de demande négatif dominant
   \item Un choc d'offre négatif dominant
   \item Un choc de productivité positif
   \item Une situation de surchauffe
}

\begin{reponse}
\textbf{(a)} La combinaison $Y \downarrow$, chômage $\uparrow$ et inflation $\downarrow$ est la \textbf{signature classique d'un choc de demande négatif} : AD se déplace vers la gauche, ce qui réduit la production et exerce une pression à la baisse sur les prix. Un choc d'offre négatif ferait aussi baisser $Y$, mais tendrait plutôt à faire \textbf{monter} l'inflation. Les prix du pétrole stables rendent en plus un choc pétrolier moins plausible.
\end{reponse}

\qcm{Pendant un ralentissement, le PIB réel recule de 2\,\%, l'investissement chute de 8\,\%, l'emploi baisse et le taux de chômage passe de 5\,\% à 7\,\%. Laquelle des affirmations suivantes décrit correctement la cyclicité observée\,?}
\choix{
   \item Le chômage est \textbf{procyclique}, comme l'investissement
   \item Le chômage est \textbf{contracyclique}, tandis que l'investissement est \textbf{procyclique}
   \item L'investissement est \textbf{contracyclique}, tandis que l'emploi est \textbf{procyclique}
   \item Le chômage et l'emploi sont tous deux \textbf{contracycliques}
}

\begin{reponse}
\textbf{(b)} Le chômage évolue en sens \textbf{inverse} du cycle : il augmente en ralentissement et diminue en expansion $\to$ il est \textbf{contracyclique}. L'investissement, au contraire, chute fortement en ralentissement et rebondit en expansion $\to$ il est \textbf{procyclique}. L'emploi est lui aussi procyclique. Les distracteurs jouent sur des inversions fréquentes entre \textit{pro} et \textit{contra}cyclique.
\end{reponse}

\qcm{Pourquoi l'investissement des entreprises tend-il à amplifier davantage les cycles que la consommation des ménages\,?}
\choix{
   \item Parce que les projets d'investissement sont souvent plus faciles à reporter quand l'incertitude augmente ou que la confiance chute
   \item Parce que l'investissement fait partie de $G$, alors que la consommation fait partie de $C$
   \item Parce que la consommation des ménages est entièrement fixe à court terme
   \item Parce que les entreprises augmentent toujours leur investissement quand la demande ralentit
}

\begin{reponse}
\textbf{(a)} L'investissement est très sensible à la \textbf{confiance}, à l'\textbf{incertitude} et aux anticipations sur la demande future. Une entreprise peut souvent \textbf{reporter} un projet d'usine ou d'équipement, alors qu'un ménage continue généralement à consommer au moins une grande partie de ses dépenses courantes. C'est pourquoi l'investissement bouge davantage que la consommation au cours du cycle. Les autres options sont fausses ou caricaturales.
\end{reponse}

\qcm{Après une baisse soudaine de la demande agrégée, pourquoi l'économie ne revient-elle pas immédiatement vers son PIB potentiel\,?}
\choix{
   \item Parce que les entreprises peuvent instantanément réduire tous les salaires nominaux
   \item Parce que les prix et salaires s'ajustent avec retard, si bien que le choc de demande se traduit d'abord par une baisse de $Y$
   \item Parce que le PIB potentiel baisse automatiquement dès que AD baisse
   \item Parce que le gouvernement compense toujours immédiatement la baisse de la demande privée
}

\begin{reponse}
\textbf{(b)} C'est précisément l'idée centrale du modèle AD-AS de court terme. Lorsque AD chute, l'économie est d'abord \textbf{poussée en dessous} de son PIB potentiel. Les entreprises font face à une baisse des ventes, mais les \textbf{prix et salaires nominaux ne s'ajustent pas instantanément}. Le retour vers le potentiel ne se fait donc pas tout de suite : à court terme, le choc se traduit d'abord par une baisse de la production et de l'emploi. C'est cette lenteur d'ajustement qui explique la pente positive de AS à court terme et qui justifie l'intervention des politiques contracycliques.
\end{reponse}

\qcm{L'écart de production est de $-3\,\%$. Cela signifie que :}
\choix{
   \item Le PIB réel est 3\,\% au-dessus du PIB potentiel
   \item Le PIB réel est 3\,\% en dessous du PIB potentiel
   \item Le PIB potentiel a diminué de 3\,\%
   \item Le taux de chômage est exactement 3\,\%
}

\begin{reponse}
\textbf{(b)} L'écart de production $= \frac{Y - Y_{\text{POT}}}{Y_{\text{POT}}} \times 100$. Un écart de $-3\,\%$ signifie que le PIB réel est 3\,\% \textbf{en dessous} de son potentiel : l'économie fonctionne sous sa capacité normale et des ressources productives sont sous-utilisées (chômage supérieur à son niveau naturel). (a) décrirait un écart \textit{positif} de $+3\,\%$ (surchauffe). (c) confond l'écart de production avec une variation du potentiel lui-même. (d) aucun lien direct entre l'écart de production et le niveau du taux de chômage.
\end{reponse}

\qcm{La courbe de demande agrégée (AD) a une pente \textbf{négative} principalement parce que :}
\choix{
   \item Les entreprises licencient quand les prix augmentent
   \item Une hausse du niveau des prix réduit la richesse réelle des ménages, qui diminuent leur consommation
   \item Le gouvernement réduit ses dépenses quand les prix montent
   \item Les importations baissent quand les prix augmentent
}

\begin{reponse}
\textbf{(b)} La pente négative de AD s'explique par l'\textbf{effet de richesse} : quand le niveau des prix ($P$) augmente, la valeur réelle des actifs des ménages (épargne, obligations) diminue $\to$ les ménages se sentent moins riches $\to$ ils réduisent leur consommation ($C \downarrow$) $\to$ la demande agrégée diminue. (a) décrit un mécanisme du côté de l'offre, pas de la demande. (c) et (d) ne sont pas les mécanismes centraux enseignés dans le modèle de base.
\end{reponse}

\qcm{La courbe d'offre agrégée (AS) a une pente \textbf{positive} à court terme parce que :}
\choix{
   \item Les consommateurs achètent plus quand les prix sont bas
   \item Les exportations augmentent quand les prix domestiques montent
   \item Les salaires nominaux sont rigides; une hausse de $P$ réduit le coût réel du travail, incitant les entreprises à produire davantage
   \item Le gouvernement augmente ses dépenses quand les prix montent
}

\begin{reponse}
\textbf{(c)} La pente positive de AS repose sur la \textbf{rigidité des salaires nominaux}. Quand le niveau des prix ($P$) augmente mais que les salaires nominaux ($w$) restent fixes (contrats, conventions), le coût réel du travail pour les entreprises diminue ($w/P \downarrow$). Les entreprises trouvent alors plus rentable de produire et d'embaucher davantage $\to$ $Y \uparrow$. D'où la relation positive entre $P$ et $Y$ le long de AS. (a) décrit la demande, pas l'offre. (b) et (d) ne sont pas les mécanismes centraux du modèle.
\end{reponse}

\qcm{Lequel des événements suivants déplace la courbe AD vers la \textbf{gauche}\,?}
\choix{
   \item Une baisse du prix des matières premières
   \item Une hausse des dépenses militaires du gouvernement
   \item Un effondrement de la confiance des consommateurs
   \item Une amélioration technologique qui réduit les coûts de production
}

\begin{reponse}
\textbf{(c)} La courbe AD se déplace lorsque les dépenses souhaitées ($C + I + G + NX$) changent \textbf{à un niveau de prix donné}. Un effondrement de la confiance réduit la consommation ($C \downarrow$) et l'investissement ($I \downarrow$) $\to$ AD vers la gauche. (a) et (d) affectent les \textbf{coûts de production} $\to$ ce sont des déplacements de AS, pas de AD. (b) augmente $G$ $\to$ AD vers la \textbf{droite}, pas la gauche.
\end{reponse}

\qcm{Lequel des événements suivants déplace la courbe AS vers la \textbf{gauche}\,?}
\choix{
   \item Une hausse de la confiance des consommateurs
   \item Une baisse des exportations due à un ralentissement mondial
   \item Des perturbations majeures des chaînes d'approvisionnement mondiales
   \item Un programme de relance budgétaire
}

\begin{reponse}
\textbf{(c)} La courbe AS se déplace quand les \textbf{coûts de production} changent à un niveau de prix donné. Des perturbations des chaînes d'approvisionnement augmentent les coûts (pénuries d'intrants, délais, surcoûts logistiques) $\to$ AS vers la gauche. (a) et (d) affectent la \textbf{demande} agrégée ($C$ et $G$), pas les coûts de production $\to$ ce sont des déplacements de AD. (b) est aussi un choc de demande ($NX \downarrow$) $\to$ AD vers la gauche.
\end{reponse}

\qcm{Un choc de demande \textbf{négatif} et un choc d'offre \textbf{négatif} ont en commun de :}
\choix{
   \item Réduire le PIB et augmenter le niveau des prix
   \item Réduire le PIB et réduire le niveau des prix
   \item Réduire le PIB, mais avoir des effets opposés sur le niveau des prix
   \item N'avoir aucun effet sur le PIB réel
}

\begin{reponse}
\textbf{(c)} Les deux types de chocs réduisent le PIB réel ($Y \downarrow$), mais leurs effets sur les prix sont \textbf{opposés} :
\begin{itemize}
   \item \textbf{Choc de demande négatif} (AD $\leftarrow$) : $Y \downarrow$ et $P \downarrow$
   \item \textbf{Choc d'offre négatif} (AS $\leftarrow$) : $Y \downarrow$ et $P \uparrow$ (stagflation)
\end{itemize}
Cette distinction est cruciale en politique économique : face à un choc de demande, la banque centrale peut stimuler l'économie sans craindre l'inflation; face à un choc d'offre, elle fait face à un dilemme.
\end{reponse}

\qcm{À la suite d'un blocage d'Ormuz, le prix du pétrole bondit et les coûts de transport augmentent fortement. Quel résultat est le plus plausible à court terme dans le modèle AD-AS\,?}
\choix{
   \item Le PIB augmente et le niveau des prix diminue
   \item Le PIB diminue et le niveau des prix diminue
   \item Le PIB diminue et le niveau des prix augmente
   \item Le PIB reste inchangé et seul le niveau des prix augmente
}

\begin{reponse}
\textbf{(c)} Un choc pétrolier est un \textbf{choc d'offre négatif} : les coûts de production augmentent dans de nombreux secteurs, si bien que AS se déplace vers la gauche. Le résultat typique est $Y \downarrow$ et $P \uparrow$ : c'est la configuration de \textbf{stagflation}. C'est aussi le type de choc qui crée un dilemme pour la banque centrale, car soutenir l'activité et combattre l'inflation exigent alors des réponses opposées.
\end{reponse}

% ============================================================================
\section{Vrai, faux ou incertain}
% ============================================================================

\setcounter{shortcounter}{0}

\vfi{Une économie peut avoir un écart de production négatif sans être en récession technique.}

\begin{reponse}
\textbf{Vrai.} Un écart de production négatif signifie simplement que $Y < Y_{\text{POT}}$ : l'économie produit en dessous de sa capacité normale. Cela peut correspondre à une récession, mais aussi à une situation de \textbf{faiblesse prolongée} où la production continue d'augmenter lentement tout en restant sous son potentiel. Une économie peut donc être \guillemotleft~molle~\guillemotright\ sans être en récession au sens usuel de baisse active de l'activité.
\end{reponse}

\vfi{Un écart de production positif est une bonne nouvelle sans ambiguïté pour l'économie à long terme.}

\begin{reponse}
\textbf{Faux.} Un écart de production positif ($Y > Y_{\text{POT}}$) signifie que l'économie produit au-delà de sa capacité normale --- c'est une situation de \textbf{surchauffe}. Cela peut sembler positif à court terme (plus d'emplois, revenus élevés), mais cela génère des pressions inflationnistes : la demande excède la capacité productive, les salaires et les prix montent. Si la surchauffe persiste, la banque centrale devra resserrer la politique monétaire, ce qui finira par provoquer un ralentissement. L'idéal est un écart de production proche de zéro, où l'économie tourne à son rythme soutenable.
\end{reponse}

\vfi{Un choc d'offre négatif est plus facile à combattre pour une banque centrale qu'un choc de demande négatif.}

\begin{reponse}
\textbf{Faux.} C'est l'inverse. Face à un choc de \textbf{demande} négatif (AD $\leftarrow$ $\to$ $Y \downarrow$, $P \downarrow$), la banque centrale peut baisser les taux pour relancer AD : cela soutient à la fois la production \textbf{et} les prix. Les deux objectifs sont \textbf{alignés}. Face à un choc d'\textbf{offre} négatif (AS $\leftarrow$ $\to$ $Y \downarrow$, $P \uparrow$), la banque centrale fait face à un \textbf{dilemme} : stimuler la demande pour soutenir $Y$ aggraverait l'inflation; resserrer pour combattre l'inflation aggraverait la récession. C'est le \guillemotleft~cauchemar~\guillemotright\ des banquiers centraux.
\end{reponse}

\vfi{Si la pandémie de 2020 avait été un choc d'offre pur (sans aucun choc de demande), on aurait observé une forte hausse de l'inflation dès le début de la crise.}

\begin{reponse}
\textbf{Vrai.} Un choc d'offre pur déplace AS vers la gauche : $Y \downarrow$ et $P \uparrow$. Si la pandémie n'avait touché que la production (fermetures d'usines, perturbation des chaînes d'approvisionnement) sans affecter la demande, les prix auraient augmenté immédiatement. Or, début 2020, l'inflation a \textbf{baissé}, ce qui montre que le choc de demande (effondrement de $C$ et $I$) a dominé dans les premiers mois. La demande a chuté plus vite que l'offre, exerçant une pression \textbf{déflationniste}. L'inflation n'est apparue qu'en 2021--2022, quand la demande a rebondi (grâce aux plans de relance) alors que l'offre restait contrainte.
\end{reponse}

\vfi{Le mécanisme d'autocorrection de l'économie fonctionne rapidement --- en quelques semaines --- grâce à l'ajustement des salaires.}

\begin{reponse}
\textbf{Faux.} Le mécanisme d'autocorrection fonctionne comme suit : en récession ($Y < Y_{\text{POT}}$), le chômage augmente $\to$ la pression salariale diminue $\to$ les coûts de production baissent $\to$ AS se déplace vers la droite $\to$ l'économie revient progressivement vers $Y_{\text{POT}}$. Mais ce processus est très \textbf{lent} : les salaires s'ajustent avec retard (environ une fois par an, et les baisses sont rares), si bien que l'autocorrection peut prendre \textbf{2 à 3 ans}. C'est précisément cette lenteur qui justifie l'intervention des politiques contracycliques (monétaire et budgétaire) pour accélérer le retour au potentiel.
\end{reponse}

\vfi{L'économie ne peut jamais produire au-delà de son PIB potentiel.}

\begin{reponse}
\textbf{Faux.} Le PIB potentiel ($Y_{\text{POT}}$) n'est pas un plafond absolu; c'est le niveau de production correspondant à l'utilisation \textbf{normale} des facteurs de production (pas maximale). L'économie peut temporairement produire au-delà de son potentiel ($Y > Y_{\text{POT}}$), en situation de \textbf{surchauffe} : les employés font des heures supplémentaires, les usines tournent au-delà de leur capacité habituelle, le chômage tombe en dessous de son taux naturel. Mais cette situation n'est pas soutenable et génère de l'inflation. C'est l'écart de production \textbf{positif}.
\end{reponse}

\vfi{Si l'on observe simultanément une baisse du PIB et une baisse des prix, on peut conclure que l'économie a subi un choc de demande négatif.}

\begin{reponse}
\textbf{Vrai.} Un choc de demande négatif déplace AD vers la gauche, ce qui produit $Y \downarrow$ et $P \downarrow$. Un choc d'offre négatif déplace AS vers la gauche, ce qui produit $Y \downarrow$ mais $P \uparrow$. L'observation conjointe de $Y \downarrow$ et $P \downarrow$ est donc la \guillemotleft~signature~\guillemotright\ d'un choc de demande. C'est exactement la logique utilisée pour interpréter la baisse de l'inflation début 2020 : elle indiquait que le choc de demande dominait. Attention cependant : si les deux types de chocs frappent simultanément, l'effet sur $P$ est ambigu et le diagnostic est plus complexe.
\end{reponse}

\vfi{Les tarifs douaniers imposés par les É.-U.\ sur les produits canadiens sont un choc de demande pur pour le Canada.}

\begin{reponse}
\textbf{Incertain.} Sans représailles canadiennes, les tarifs américains réduisent les exportations du Canada ($NX \downarrow$), ce qui est un choc de \textbf{demande} négatif (AD $\leftarrow$). Mais \textbf{avec} des représailles (tarifs canadiens sur les produits américains), les intrants importés coûtent plus cher pour les entreprises canadiennes, ce qui est un choc d'\textbf{offre} négatif (AS $\leftarrow$). Le résultat dépend donc de la politique de représailles : sans représailles, c'est principalement un choc de demande ($Y \downarrow$, $P \downarrow$); avec représailles, c'est un double choc ($Y \downarrow\downarrow$, $P$ ambigu, risque de stagflation). C'est exactement la distinction faite en cours.
\end{reponse}

% ============================================================================
\section{Questions à développement}
% ============================================================================

\setcounter{shortcounter}{0}

\shortq{Considérez une économie initialement au plein emploi ($Y = Y_{\text{POT}}$). Une sécheresse sévère détruit une partie importante des récoltes et perturbe la production alimentaire du pays.}

\begin{enumerate}[label=(\alph*)]
   \item S'agit-il d'un choc de demande ou d'un choc d'offre\,? Justifiez.
   \item Dans quel sens la courbe concernée se déplace-t-elle\,? Décrivez qualitativement le nouvel équilibre de court terme (effet sur $Y$ et $P$).
   \item Comment appelle-t-on la combinaison de stagnation économique et d'inflation\,?
   \item Expliquez pourquoi la banque centrale fait face à un dilemme face à ce type de choc.
\end{enumerate}

\begin{reponse}
\textbf{(a)} C'est un \textbf{choc d'offre négatif}. La sécheresse détruit des récoltes et perturbe la production $\to$ les coûts de production augmentent (pénurie d'intrants alimentaires, hausse des prix agricoles) et la capacité productive diminue. La demande des ménages n'est pas directement touchée (ils veulent toujours consommer autant), mais les entreprises ne peuvent plus produire autant à chaque niveau de prix.

\textbf{(b)} La courbe \textbf{AS se déplace vers la gauche}. Le nouvel équilibre de court terme se situe à l'intersection de AD (inchangée) et AS$'$ (déplacée vers la gauche). Le résultat est : $Y \downarrow$ (la production diminue) et $P \uparrow$ (les prix augmentent, car les coûts de production sont plus élevés).

\textbf{(c)} C'est la \textbf{stagflation} : la contraction économique (stagnation, voire récession) coexiste avec une hausse des prix (inflation).

\textbf{(d)} La banque centrale fait face à un \textbf{dilemme} car ses deux objectifs sont en conflit :
\begin{itemize}
   \item \textbf{Combattre l'inflation} : resserrer la politique monétaire (hausse des taux) $\to$ AD vers la gauche $\to$ l'inflation diminue, mais la récession \textbf{s'aggrave}
   \item \textbf{Soutenir la production} : assouplir la politique monétaire (baisse des taux) $\to$ AD vers la droite $\to$ la production remonte, mais l'inflation \textbf{s'aggrave}
\end{itemize}
Contrairement à un choc de demande (où les deux objectifs sont alignés), un choc d'offre oblige à choisir entre deux maux.
\end{reponse}

\shortq{Le gouvernement d'un petit pays annonce une hausse massive des dépenses en infrastructure ($G \uparrow$) financée par l'emprunt.}

\begin{enumerate}[label=(\alph*)]
   \item S'agit-il d'un choc de demande ou d'un choc d'offre\,? Quelle courbe se déplace et dans quelle direction\,?
   \item Si l'économie est initialement en récession ($Y < Y_{\text{POT}}$), illustrez qualitativement l'effet à court terme sur $Y$ et $P$.
   \item Si l'économie est initialement en surchauffe ($Y > Y_{\text{POT}}$), quel serait le principal risque associé à cette même politique\,?
   \item Supposons maintenant que l'investissement en infrastructure augmente durablement la productivité ($A \uparrow$). Quelle courbe supplémentaire se déplace\,? Comparez l'effet global sur $Y$ et $P$ avec la situation en (b).
\end{enumerate}

\begin{reponse}
\textbf{(a)} C'est un \textbf{choc de demande positif}. La hausse de $G$ augmente directement les dépenses souhaitées à chaque niveau de prix $\to$ la courbe \textbf{AD se déplace vers la droite}.

\textbf{(b)} Partant d'une récession ($Y < Y_{\text{POT}}$), le déplacement de AD vers la droite pousse l'équilibre le long de la courbe AS (croissante) : $Y \uparrow$ et $P \uparrow$. Si le déplacement est bien calibré, l'économie se rapproche de $Y_{\text{POT}}$, ce qui réduit l'écart de production et le chômage. La hausse des prix est modérée car l'économie avait des capacités inutilisées.

\textbf{(c)} Si l'économie est déjà en surchauffe ($Y > Y_{\text{POT}}$), une relance budgétaire supplémentaire déplacerait AD encore plus vers la droite. L'économie opérant déjà au-delà de ses capacités normales, l'essentiel du stimulus se traduirait par de l'\textbf{inflation} plutôt que par des gains de production. La hausse de $P$ serait forte, la hausse de $Y$ modeste. C'est une politique \textbf{inappropriée} en contexte de surchauffe.

\textbf{(d)} Si l'infrastructure augmente la productivité ($A \uparrow$), les courbes \textbf{AS se déplace aussi vers la droite} (les entreprises produisent davantage à chaque niveau de prix grâce à la meilleure productivité). De plus, le PIB potentiel $Y_{\text{POT}}$ augmente. L'effet combiné (AD $\rightarrow$ et AS $\rightarrow$) produit une hausse de $Y$ plus \textbf{importante} qu'en (b), avec un effet sur $P$ \textbf{ambigu} (les deux déplacements se renforcent pour $Y$ mais s'opposent pour $P$). C'est pourquoi les investissements publics productifs (qui augmentent $A$) sont économiquement plus désirables que les dépenses courantes.
\end{reponse}

\shortq{Complétez le tableau suivant pour deux épisodes vus en séance : (i) des tarifs américains sur les exportations canadiennes sans représailles canadiennes, (ii) un choc pétrolier lié à un blocage d'Ormuz. Justifiez brièvement chaque ligne.}

\begin{enumerate}[label=(\alph*)]
   \item Pour chaque épisode, indiquez : le choc dominant, le déplacement dans AD-AS, l'effet attendu sur $Y$, l'effet attendu sur $P$, et si le dilemme de la banque centrale est faible, moyen ou fort.
   \item Reprenez maintenant les deux épisodes en ajoutant un canal supplémentaire dans chaque cas : (i) en réponse aux tarifs américains, le Canada impose des représailles sur plusieurs importations américaines; (ii) après le choc d'Ormuz, la hausse du prix du pétrole augmente les profits et l'investissement du secteur pétrolier canadien. Pour chaque épisode, indiquez quelle courbe additionnelle se déplace, dans quel sens, et comment cela modifie le diagnostic sur $Y$, $P$ et le dilemme de la banque centrale.
\end{enumerate}

\begin{reponse}
\textbf{(a)} Un bon tableau de réponse est le suivant :

\begin{center}
\renewcommand{\arraystretch}{1.2}
\begin{tabular}{p{3.2cm} p{2.5cm} p{2.2cm} p{1.5cm} p{1.5cm} p{2.4cm}}
\toprule
\textbf{Épisode} & \textbf{Choc dominant} & \textbf{Déplacement} & \textbf{$Y$} & \textbf{$P$} & \textbf{Dilemme BC} \\
\midrule
Tarifs sans représailles & Demande négative & AD $\leftarrow$ & $\downarrow$ & $\downarrow$ & Faible \\
Blocage d'Ormuz & Offre négative & AS $\leftarrow$ & $\downarrow$ & $\uparrow$ & Fort \\
\bottomrule
\end{tabular}
\end{center}

Justification :
\begin{itemize}
   \item \textbf{Tarifs sans représailles :} les exportations canadiennes diminuent ($NX \downarrow$), donc AD se déplace vers la gauche. C'est le cas le plus proche d'un choc de demande standard.
   \item \textbf{Blocage d'Ormuz :} les coûts énergétiques et de transport augmentent dans l'ensemble de l'économie, ce qui déplace AS vers la gauche et crée une situation de stagflation.
\end{itemize}

\textbf{(b)} Ici, chaque épisode devient un \textbf{choc mixte}.
\begin{itemize}
   \item \textbf{Tarifs avec représailles canadiennes :} au choc initial de \textbf{demande négative} (AD $\leftarrow$ car $NX \downarrow$) s'ajoute un choc d'\textbf{offre négative} (AS $\leftarrow$) parce que certains intrants importés des États-Unis deviennent plus coûteux. Le recul de $Y$ devient \textbf{plus fort} qu'en (a). L'effet sur $P$ devient \textbf{ambigu ou légèrement haussier} : la demande pousse les prix vers le bas, mais la hausse des coûts pousse les prix vers le haut. Le dilemme de la banque centrale devient donc \textbf{plus important}.
   \item \textbf{Choc d'Ormuz avec profits pétroliers canadiens :} au choc initial d'\textbf{offre négative} (AS $\leftarrow$) s'ajoute un choc de \textbf{demande positive} (AD $\rightarrow$) via les profits, l'investissement et potentiellement les revenus dans le secteur énergétique canadien. L'effet sur $P$ devient \textbf{encore plus inflationniste} car les deux canaux poussent les prix vers le haut. L'effet sur $Y$ devient \textbf{ambigu ou moins négatif} : AS tire $Y$ vers le bas, mais AD soutient l'activité dans les régions productrices. Le dilemme de la banque centrale devient \textbf{très fort}.
\end{itemize}
Autrement dit, ces complications rapprochent davantage les deux épisodes de cas réels : les chocs ne déplacent pas toujours une seule courbe, et le diagnostic monétaire devient plus difficile lorsque plusieurs mécanismes agissent en sens opposés sur $Y$ ou dans le même sens sur $P$.
\end{reponse}

\end{document}
